\chapter{Materials and Methods}

\section{System Flow Diagram of the Computational Framework}

\subsection{Overview of the Computational Framework}

The overall methodology follows a systematic five-phase pipeline as illustrated in
Figure~\ref{fig:workflow}. All code was implemented in Python 3.14 using
exclusively open-source libraries and executed on a macOS system (CPU).

\begin{figure}[htbp]
\centering
\begin{tikzpicture}[
  node distance=0.55cm,
  phase/.style={rectangle, rounded corners=6pt, draw=black, thick,
                fill=blue!20, minimum width=7.6cm, minimum height=0.7cm,
                align=center, font=\bfseries\small},
  proc/.style={rectangle, rounded corners=4pt, draw=black,
               fill=cyan!12, minimum width=7.2cm, minimum height=0.5cm,
               align=center, font=\small},
  final/.style={rectangle, rounded corners=6pt, draw=black, very thick,
                fill=green!25, minimum width=7.6cm, minimum height=0.7cm,
                align=center, font=\bfseries\small},
  arr/.style={->, thick, >=latex}
]
\node[phase] (ph1) {Phase 1: Data Acquisition \& Preprocessing};
\node[proc, below=0.2cm of ph1] (p1a) {ChEMBL CHEMBL203 retrieval via REST API};
\node[proc, below=of p1a] (p1b) {pIC$_{50}$ conversion, binary labelling, RDKit curation};
\node[proc, below=of p1b] (p1c) {Exploratory data analysis (class balance, scaffold diversity)};
\node[phase, below=0.5cm of p1c] (ph2) {Phase 2: Molecular Graph Construction};
\node[proc, below=0.2cm of ph2] (p2a) {29-dim node features + 12-dim edge features};
\node[phase, below=0.5cm of p2a] (ph3) {Phase 3: Model Development \& Training};
\node[proc, below=0.2cm of ph3] (p3a) {\textbf{5 GNNs} (GCN, GAT, GIN, AttentiveFP, GraphSAGE)\\\textbf{vs.\ 2 traditional ML baselines} (Random Forest, MLP)};
\node[phase, below=0.5cm of p3a] (ph4) {Phase 4: Evaluation \& Benchmarking};
\node[proc, below=0.2cm of ph4] (p4a) {AUROC, AUPRC, F1, MCC on scaffold test set};
\node[phase, below=0.5cm of p4a] (ph5) {Phase 5: Nepali Plant Virtual Screening};
\node[proc, below=0.2cm of ph5] (p5a) {\textbf{Select 7 Nepali medicinal plant species}\\(documented anti-cancer / anti-inflammatory use $+$ COCONUT~2.0 coverage)};
\node[proc, below=of p5a] (p5b) {COCONUT 2.0 retrieval + RDKit standardisation $\rightarrow$ 535 compounds};
\node[proc, below=of p5b] (p5c) {Score with \textbf{best models} (GCN $+$ Random Forest consensus)\\$+$ Lipinski Ro5, QED, ADMET-proxy filtering};
\node[final, below=0.5cm of p5c] (out) {Ranked Nepali Natural-Product Lead Candidates};

\foreach \a/\b in {ph1/p1a, p1a/p1b, p1b/p1c, p1c/ph2, ph2/p2a, p2a/ph3,
                   ph3/p3a, p3a/ph4, ph4/p4a, p4a/ph5,
                   ph5/p5a, p5a/p5b, p5b/p5c, p5c/out}
  \draw[arr] (\a) -- (\b);
\end{tikzpicture}
\caption{Five-phase computational pipeline: benchmarking five GNN architectures
         against traditional machine-learning baselines for EGFR inhibitor
         classification, and deploying the best-performing models (GCN and Random
         Forest) to screen a Nepali medicinal-plant compound library---whose seven
         species were selected by documented anti-cancer use and COCONUT~2.0 coverage.}
\label{fig:workflow}
\end{figure}

\section{Data Acquisition}

\subsection{Target Selection}

The primary target for bioactivity data retrieval was the Epidermal Growth Factor
Receptor (EGFR), indexed in ChEMBL as target identifier \texttt{CHEMBL203}. EGFR
was selected as the model target because it is a clinically validated oncology
target in NSCLC, and CHEMBL203 aggregates the largest publicly available collection
of quantitative EGFR inhibition bioactivity measurements.

\subsection{Bioactivity Data Retrieval}

Bioactivity data was retrieved programmatically from the ChEMBL database using the
\texttt{chembl-webresource-client} Python library, which queries the then-current
ChEMBL release through the REST API (retrieved May 2026). The query was
restricted to compounds satisfying the following criteria:
\begin{itemize}
    \item \textbf{Standard type}: IC$_{50}$ (half-maximal inhibitory concentration)
    \item \textbf{Standard relation}: Exact value only (\texttt{=})
\end{itemize}

Records with missing SMILES strings, missing IC$_{50}$ values, or non-exact
relations ($>$, $<$, $\geq$, $\leq$) were excluded. IC$_{50}$ values reported
in ChEMBL for this target are in nanomolar (nM) units, as assumed throughout
the pIC$_{50}$ conversion. The raw dataset contained
approximately 19,000 compound-activity pairs prior to deduplication and cleaning.

\subsection{Data Cleaning and Labelling}

The raw IC$_{50}$ values (in nanomolar units) were converted to the logarithmic
pIC$_{50}$ scale by taking the negative base-ten logarithm of the molar
concentration, so that more potent compounds receive higher pIC$_{50}$ scores.
Binary activity labels were then assigned using the standard medicinal chemistry
threshold: compounds with a pIC$_{50}$ of at least 6.0 were labelled \textit{active},
and those below 6.0 were labelled \textit{inactive}. This threshold corresponds to
an IC$_{50}$ of $\leq 1~\mu$M. For compounds with multiple reported measurements,
the median pIC$_{50}$ was used to obtain a single robust value per molecule.

\subsection{Molecular Standardisation and Filtering}

Molecular standardisation was performed using RDKit (version 2026.3.2):
\begin{itemize}
    \item \textbf{Salt removal}: The largest fragment was retained for each compound.
    \item \textbf{Canonicalisation}: SMILES strings were canonicalised to ensure
          unique molecular representation.
    \item \textbf{Molecular weight filter}: Compounds with MW $>$ 1000 Da
          were excluded.
    \item \textbf{Deduplication}: Duplicate canonical SMILES were collapsed using
          the median pIC$_{50}$.
\end{itemize}

The final curated dataset comprised \textbf{10,523 unique compounds}, distributed
as 7,792 actives (74.0\%) and 2,731 inactives (26.0\%).

\subsection{Exploratory Data Analysis}

Prior to molecular graph construction, an exploratory data analysis (EDA) step was
performed on the cleaned dataset to characterise its statistical structure and inform
modelling decisions. Using RDKit and pandas, the analysis quantified the class
balance, the distribution of pIC$_{50}$ potency values, the ranges of the
physicochemical properties (molecular weight, LogP), and the Bemis-Murcko scaffold
diversity (number of unique scaffolds and the proportion of singleton scaffolds).
This characterisation, reported in Chapter~4, both confirms the integrity of the
curated dataset and motivates the use of scaffold-based splitting, given the
long-tailed scaffold distribution observed.

\section{Molecular Graph Construction}

\subsection{Graph Representation}

Each molecule was converted to an undirected graph using
PyTorch Geometric (version 2.7.0), where every node represents an atom and every
edge represents a chemical bond (the molecule-to-graph conversion is illustrated in
Figure~\ref{fig:intro_mol2graph}). Each atom node carries a 29-dimensional feature
vector and each bond edge a 12-dimensional feature vector, detailed below.

\subsection{Node Features}

Each atom was represented by a 29-dimensional feature vector encoding:
\begin{itemize}
    \item \textbf{Atom type} (one-hot, 10 dim): C, N, O, S, F, Cl, Br, I, P, Unknown
    \item \textbf{Atom degree} (one-hot, 6 dim): 0, 1, 2, 3, 4, $\geq$5
    \item \textbf{Formal charge} (1 dim): integer value
    \item \textbf{Hydrogen count} (one-hot, 5 dim): 0, 1, 2, 3, $\geq$4
    \item \textbf{Hybridisation} (one-hot, 5 dim): SP, SP2, SP3, SP3D, SP3D2
    \item \textbf{Aromaticity} (binary, 1 dim)
    \item \textbf{Ring membership} (binary, 1 dim)
\end{itemize}

\subsection{Edge Features}

Each bond was represented by a 12-dimensional edge feature vector encoding:
\begin{itemize}
    \item \textbf{Bond type} (one-hot, 4 dim): Single, Double, Triple, Aromatic
    \item \textbf{Conjugation} (binary, 1 dim)
    \item \textbf{Ring membership} (binary, 1 dim)
    \item \textbf{Stereochemistry} (one-hot, 6 dim): STEREONONE, STEREOANY,
          STEREOZ, STEREOE, STEREOCIS, STEREOTRANS
\end{itemize}

\section{Data Splitting}

To ensure rigorous evaluation on genuinely novel chemical scaffolds, a deterministic
\textbf{Bemis-Murcko Scaffold Split} was implemented using RDKit. The Bemis-Murcko
algorithm extracts the core macrocyclic framework of each molecule by retaining ring
systems, linker atoms, and chain atoms connected to rings while removing side chains.
Compounds sharing identical scaffold hashes were allocated exclusively to one split
partition. The dataset was partitioned as follows:

\begin{table}[H]
\centering
\caption{Dataset partitioning using the Bemis-Murcko scaffold split.}
\label{tab:split}
\begin{tabular}{lccc}
\toprule
\textbf{Partition} & \textbf{Proportion} & \textbf{Compounds} & \textbf{Unique Scaffolds} \\
\midrule
Training   & 80\% & 8,418  & majority \\
Validation & 10\% & 1,052  & held-out \\
Test       & 10\% & 1,053  & held-out \\
\bottomrule
\end{tabular}
\end{table}

\section{Machine Learning Model Development}

\subsection{Traditional Baseline Models}

Two traditional machine learning baselines were implemented using
scikit-learn (version 1.8.0) trained on 2048-bit ECFP4 (Morgan) fingerprints
(radius = 2) computed by RDKit:

\begin{itemize}
    \item \textbf{Random Forest (RF)}: An ensemble of 500 decision trees with
          \texttt{class\_weight=} \texttt{'balanced'} to account for the 74:26 class imbalance.
          Hyperparameters included \texttt{max\_features='sqrt'} and
          \texttt{min\_samples\_leaf=1}.
    \item \textbf{Multi-Layer Perceptron (MLP)}: A three-layer neural network
          (2048$\to$512$\to$128$\to$1) with BatchNorm and Dropout(0.3) at each hidden
          layer, trained on ECFP4 fingerprints with BCEWithLogitsLoss (pos\_weight),
          Adam optimiser, and early stopping (patience = 10) on an internal 90/10
          validation split. This evaluates whether deep learning on fingerprint features
          alone offers benefits beyond the Random Forest ensemble.
\end{itemize}

\subsection{Graph Neural Network Architectures}

Five GNN architectures were implemented using PyTorch Geometric:

\subsubsection{Graph Convolutional Network (GCN)}

The GCN propagates information using degree-normalised spectral convolution. At
each layer, every atom embedding is updated by averaging the transformed
embeddings of its bonded neighbours and itself, where the averaging weights are
normalised by the degrees of the participating atoms (with added self-loops). This
symmetric normalisation prevents high-degree atoms from dominating the
representation. Three convolutional layers were stacked, each followed by batch
normalisation, a ReLU non-linearity, and dropout.

\subsubsection{Graph Attention Network (GAT)}

The GAT replaces the fixed normalisation of the GCN with learnable attention
weights: for each atom, the contribution of every neighbour is weighted by an
attention coefficient computed from the two atoms' embeddings and the connecting
bond features, then normalised across all neighbours. This lets the network learn
which bonds and substructures are most relevant to activity. Four attention heads
were used in the intermediate layers (their outputs concatenated) and a single
averaged head in the final layer. Bond (edge) features were supplied directly to
the attention mechanism through the GATv2 formulation.

\subsubsection{Graph Isomorphism Network (GIN)}

GIN is the most theoretically expressive message-passing scheme within the
Weisfeiler--Lehman framework. Rather than averaging, each atom embedding is updated
by summing its neighbours' embeddings (preserving multiset information that mean or
max aggregation would discard), adding the atom's own embedding, and passing the
result through a two-layer multilayer perceptron. The GINEConv variant used here
extends this update to incorporate the 12-dimensional bond feature vectors.

\subsubsection{GraphSAGE}

GraphSAGE \parencite{hamilton2017graphsage} updates each atom by combining its own
previous embedding with the mean of its neighbours' embeddings and passing the
combination through a learnable transformation. Because the update depends only on
local neighbourhood aggregation rather than the global graph structure, GraphSAGE
is inherently inductive and generalises naturally to unseen molecules. Three
SAGEConv layers with batch normalisation and ReLU activation are used, followed by
concatenated global mean and max pooling and a linear classification head.
GraphSAGE was included to provide an inductive neighbourhood-sampling baseline that
is evaluated under the same scaffold split conditions as the attention-based models.

\subsubsection{AttentiveFP}

AttentiveFP implements a two-level attention mechanism. At the atom level, a graph
attention layer aggregates neighbourhood information weighted by learned
atom--atom attention scores. At the molecule level, a Gated Recurrent Unit
(GRU)-based readout iteratively refines a single global molecular embedding over
several timesteps, allowing the model to focus on the substructures most predictive
of activity. Crucially, AttentiveFP natively accepts both the atom and bond feature
vectors without preprocessing, making it uniquely suited to exploit the full
29-dimensional node and 12-dimensional edge feature specification.

\subsection{Training Protocol}

All GNN models were trained with the following protocol:
\begin{itemize}
    \item \textbf{Loss function}: Binary cross-entropy with logits incorporating
          a positive class weight (the ratio of negative to positive training
          examples, approximately 0.34) to address the class imbalance.
    \item \textbf{Optimiser}: Adam with its default momentum parameters.
    \item \textbf{Scheduler}: \texttt{ReduceLROnPlateau} with reduction factor 0.5
          and patience 10 epochs.
    \item \textbf{Early stopping}: Training halted if validation AUROC did not
          improve for 20 consecutive epochs (30 for the tuned AttentiveFP).
    \item \textbf{Batch size}: 64 for default training; selected by Optuna for the
          tuned AttentiveFP.
    \item \textbf{Readout}: Global mean and max pooling were concatenated to form
          the graph-level representation.
\end{itemize}

\subsection{Model Evaluation Metrics}

Models were evaluated using the following metrics on the scaffold-held-out test set:

\begin{itemize}
    \item \textbf{AUROC}: Threshold-independent measure of discriminative ability
          across the full range of classification thresholds.
    \item \textbf{AUPRC}: Area under the precision-recall curve, more informative
          than AUROC under class imbalance.
    \item \textbf{F1-Score}: Harmonic mean of precision and recall at the 0.5
          decision threshold.
    \item \textbf{Matthews Correlation Coefficient (MCC)}: A balanced correlation
          between predicted and true labels derived from all four confusion-matrix
          counts (true/false positives and negatives), ranging from $-1$ (total
          disagreement) through $0$ (random) to $+1$ (perfect prediction). It is
          robust to class imbalance and is reported as the primary balanced metric.
    \item \textbf{Accuracy}: Proportion of correctly classified compounds.
\end{itemize}

\subsection{Hyperparameter Tuning}

Bayesian hyperparameter optimisation was performed for the AttentiveFP model using
the Optuna framework (version 4.8.0) with the Tree-structured Parzen Estimator (TPE)
sampler. The search space encompassed:

\begin{table}[H]
\centering
\caption{Hyperparameter search space for AttentiveFP Optuna optimisation.}
\label{tab:search}
\begin{tabular}{lll}
\toprule
\textbf{Hyperparameter} & \textbf{Type} & \textbf{Range / Options} \\
\midrule
\texttt{hidden\_channels}  & Categorical   & \{64, 128, 256\}           \\
\texttt{num\_layers}       & Integer       & 2--4                        \\
\texttt{num\_timesteps}    & Integer       & 2--4                        \\
\texttt{dropout}           & Float (log)   & 0.1--0.5                    \\
\texttt{batch\_size}       & Categorical   & \{32, 64, 128\}             \\
\texttt{learning\_rate}    & Float (log)   & $10^{-4}$--$10^{-2}$        \\
\texttt{weight\_decay}     & Float (log)   & $10^{-5}$--$10^{-3}$        \\
\bottomrule
\end{tabular}
\end{table}

The objective function to maximise was validation AUROC. Pruning of unpromising
trials was performed using the Optuna \texttt{MedianPruner}.

\section{Nepali Medicinal Plant Virtual Screening}

\subsection{Plant Species and Compound Library Curation}

A virtual screening library was assembled from the COCONUT 2.0 natural product
database \parencite{sorokina2021coconut} by filtering for compounds annotated to
seven traditional Nepali medicinal plant species:

\begin{table}[H]
\centering
\caption{Nepali medicinal plant species used for virtual screening.}
\label{tab:plants}
\begin{tabular}{llc}
\toprule
\textbf{Species} & \textbf{Common Name} & \textbf{Traditional Use} \\
\midrule
\textit{Swertia chirayita}     & Chirayito       & Fever, liver conditions \\
\textit{Rhododendron arboreum} & Lali Gurans     & Anti-inflammatory        \\
\textit{Nardostachys jatamansi}& Jatamansi       & Neurological, cardiac    \\
\textit{Terminalia chebula}    & Haritaki        & Antioxidant, Triphala    \\
\textit{Berberis aristata}     & Chutro          & Anticancer (berberine)   \\
\textit{Ocimum sanctum}        & Tulsi           & Anticancer, anti-EGFR    \\
\textit{Tinospora cordifolia}  & Giloy/Gurjo     & Immunomodulatory         \\
\bottomrule
\end{tabular}
\end{table}

Retrieved compounds were subjected to the same preprocessing pipeline as the
ChEMBL training data: salt removal (largest fragment retained), SMILES
canonicalisation, molecular weight filter (MW $\leq$ 1000 Da), removal of
invalid SMILES, and deduplication by canonical SMILES. ADMET-proxy descriptors
(MW, LogP, HBD, HBA, TPSA, QED) were computed using RDKit for all retained
compounds.

\subsection{Inference and Prioritisation}

To avoid basing lead selection on any single model---and in particular to avoid the
weakest classifier---the library was scored with the study's two best-performing
models: the best graph neural network (GCN) and the best overall classifier (Random
Forest). For the GCN, each compound was featurised using the identical 29-dimensional
node and 12-dimensional edge graph specification as the training data and passed
through a forward inference to produce $\hat{p}_i = P(\text{active} \mid G_i)$. For
the Random Forest, each compound was encoded as a 2048-bit ECFP4 fingerprint and
scored with the trained ensemble. Predictions from both models were retained so that
consensus (agreement between the two best models) could be assessed. The tuned
AttentiveFP screen was also computed for comparison, exposing its over-confidence
relative to the two best-calibrated models.

Compounds were ranked by predicted activity probability $\hat{p}_i$ under each model.
The top leads were further assessed by:
\begin{itemize}
    \item \textbf{Lipinski Rule of Five compliance}: MW $\leq$ 500 Da, LogP $\leq$ 5,
          HBD $\leq$ 5, HBA $\leq$ 10 (necessary condition for oral bioavailability).
    \item \textbf{QED score}: A composite drug-likeness score targeting QED $\geq$ 0.40.
    \item \textbf{TPSA}: Values $\leq$ 140 \AA$^2$ for adequate oral absorption.
\end{itemize}

The final outputs are prioritised ranked lists for each best model
(\texttt{results/nepali\_leads\_rf.csv} and \texttt{results/nepali\_leads\_gcn.csv})
of Nepali natural-product-derived compounds with computed ADMET-proxy properties,
representing actionable candidates for experimental EGFR inhibition assays.

